\section{Análisis de la Solución Propuesta por el Operador}

En esta sección se analiza la solución propuesta por el operador para el escenario determinista. Esta solución considera las siguientes rutas:

\begin{equation*}
R_{T1} = [1,3], \qquad R_{T2} = [2,4]
\end{equation*}

Es decir, el camión $T1$ realiza la ruta Depósito->Estación 1->Estación 3->Depósito, mientras que el camión $T2$ recorre la ruta Depósito->Estación 2->Estación 4->Depósito. El operador declara un costo de distancia igual a \$260 y shortage igual a cero. A continuación se verifica la factibilidad de dicha solución.

\subsection{Verificación de Capacidad de Compartimentos}

Para el camión $T1$, que visita las estaciones 1 y 3, las demandas totales son:

\begin{align*}
Regular_{T1} &= 3000 + 2500 = 5500 \text{ L} \\
Diesel_{T1} &= 2000 + 3000 = 5000 \text{ L}
\end{align*}

Considerando que en $T1$ el compartimento $C0$ transporta Regular y $C1$ transporta Diésel, se obtiene:

\begin{align*}
C0_{T1} &: 5500 \leq 8000 \\
C1_{T1} &: 5000 \leq 7000
\end{align*}

Por lo tanto, el camión $T1$ cumple las restricciones de capacidad.

Para el camión $T2$, que visita las estaciones 2 y 4, las demandas totales son:

\begin{align*}
Regular_{T2} &= 4000 + 1000 = 5000 \text{ L} \\
Diesel_{T2} &= 1500 + 2500 = 4000 \text{ L}
\end{align*}

Considerando que en $T2$ el compartimento $C0$ transporta Diésel y $C1$ transporta Regular, se obtiene:

\begin{align*}
C0_{T2} &: 4000 \leq 6000 \\
C1_{T2} &: 5000 \leq 9000
\end{align*}

Por lo tanto, el camión $T2$ también cumple las restricciones de capacidad.

\begin{table}[H]
\centering
\caption{Verificación de capacidad de la solución del operador}
\begin{tabular}{ccccc}
\toprule
Camión & Compartimento & Producto & Carga requerida (L) & Capacidad (L) \\
\midrule
T1 & C0 & Regular & 5500 & 8000 \\
T1 & C1 & Diésel & 5000 & 7000 \\
T2 & C0 & Diésel & 4000 & 6000 \\
T2 & C1 & Regular & 5000 & 9000 \\
\bottomrule
\end{tabular}
\end{table}

Aunque las capacidades se cumplen satisfactoriamente, es posible observar que los \textit{fill ratios} declarados por el operador no coinciden realmente con las cargas requeridas por las rutas. Por ejemplo, para $T1$, el fill ratio para $C0$ debería ser $5500/8000 = 68.75\%$, no $62.5\%$.

\subsection{Verificación de Estabilidad de Carga}

La estabilidad de carga exige que la diferencia entre los \textit{fill ratios} de los compartimentos de un mismo camión no supere $\Delta = 0.30$ en ningún instante de la ruta.

Para el camión $T1$, los fill ratios son:

\begin{table}[H]
\centering
\caption{Fill ratios del camión T1}
\begin{tabular}{cccc}
\toprule
Instante & Fill ratio C0 & Fill ratio C1 & Diferencia \\
\midrule
Salida depósito & $5500/8000 = 0.6875$ & $5000/7000 = 0.7143$ & 0.0268 \\
Después Est. 1 & $2500/8000 = 0.3125$ & $3000/7000 = 0.4286$ & 0.1161 \\
Después Est. 3 & $0/8000 = 0$ & $0/7000 = 0$ & 0 \\
\bottomrule
\end{tabular}
\end{table}

Por lo tanto, el camión $T1$ cumple la restricción de estabilidad.

Para el camión $T2$, los fill ratios son:

\begin{table}[H]
\centering
\caption{Fill ratios del camión T2}
\begin{tabular}{cccc}
\toprule
Instante & Fill ratio C0 & Fill ratio C1 & Diferencia \\
\midrule
Salida depósito & $4000/6000 = 0.6667$ & $5000/9000 = 0.5556$ & 0.1111 \\
Después Est. 2 & $2500/6000 = 0.4167$ & $1000/9000 = 0.1111$ & 0.3056 \\
Después Est. 4 & $0/6000 = 0$ & $0/9000 = 0$ & 0 \\
\bottomrule
\end{tabular}
\end{table}

En el caso de $T2$, después de atender la estación 2, la diferencia entre fill ratios es:

\begin{equation*}
|0.4167 - 0.1111| = 0.3056
\end{equation*}

Dado que $0.3056 > 0.30$, la solución propuesta por el operador incumple de manera leve la restricción de estabilidad de carga para el camión $T2$.

\subsection{Verificación de Ventanas de Tiempo}


La velocidad de los camiones es de 60 km/h, por lo que cada kilómetro equivale aproximadamente a un minuto de viaje. Además, el tiempo de servicio en cada estación es de 30 minutos.

Para el camión $T1$, cuya hora de salida es 05:00, se obtiene:

\begin{table}[H]
\centering
\caption{Verificación de ventanas de tiempo para T1}
\begin{tabular}{ccccc}
\toprule
Tramo & Llegada & Inicio servicio & Ventana & Estado \\
\midrule
Depósito--Est. 1 & 05:20 & 06:00 & [06:00--10:00] & Cumple con espera \\
Est. 1--Est. 3 & 06:55 & 08:00 & [08:00--14:00] & Cumple con espera \\
\bottomrule
\end{tabular}
\end{table}

El cálculo anterior es obtenido considerando que el camión $T1$ sale a las 05:00 y demora 20 minutos en llegar a la estación 1. Como llega a las 05:20 y la ventana comienza a las 06:00, debe esperar unos 40 minutos. Luego realiza su servicio durante 30 minutos, saliendo a las 06:30. Desde la estación 1 hasta la estación 3 recorre 25 km, por lo que llega a las 06:55. Como la estación 3 comienza a recibir a las 08:00, el camión espera hasta dicha hora.

Para el camión $T2$, cuya hora de salida es 05:30, se obtiene:

\begin{table}[H]
\centering
\caption{Verificación de ventanas de tiempo para T2}
\begin{tabular}{ccccc}
\toprule
Tramo & Llegada & Inicio servicio & Ventana & Estado \\
\midrule
Depósito--Est. 2 & 06:05 & 07:00 & [07:00--12:00] & Cumple con espera \\
Est. 2--Est. 4 & 07:45 & 07:45 & [06:00--11:00] & Cumple \\
\bottomrule
\end{tabular}
\end{table}

En este caso, el camión $T2$ sale a las 05:30 y se demora unos 35 minutos en llegar a la estación 2, llegando a las 06:05. Como la ventana comienza a las 07:00, espera hasta esa hora y luego realiza el servicio durante 30 minutos, saliendo a las 07:30. Posteriormente recorre 15 km hacia la estación 4, llegando a las 07:45, lo cual se encuentra dentro de la ventana [06:00--11:00].

Por lo tanto, considerando espera anticipada, la solución del operador cumple las ventanas de tiempo.

\subsection{Verificación del Costo Declarado}

El operador declara un costo de distancia igual a \$260. Sin embargo, al calcular la distancia real recorrida se obtiene:

\begin{align*}
Distancia_{T1} &= d_{0,1} + d_{1,3} + d_{3,0} \\
&= 20 + 25 + 15 = 60 \text{ km}
\end{align*}

\begin{align*}
Distancia_{T2} &= d_{0,2} + d_{2,4} + d_{4,0} \\
&= 35 + 15 + 40 = 90 \text{ km}
\end{align*}

Por lo tanto, la distancia total recorrida es:

\begin{equation*}
Distancia_{total} = 60 + 90 = 150 \text{ km}
\end{equation*}

Como el costo por kilómetro es de \$2/km, el costo real por distancia es:

\begin{equation*}
C_{dist} = 150 \cdot 2 = 300
\end{equation*}

Además, ambos camiones son utilizados, por lo que se deben considerar los costos fijos:

\begin{equation*}
C_{fijo} = 500 + 400 = 900
\end{equation*}

Si se considera que no existe shortage, la penalización por demanda no satisfecha es cero:

\begin{equation*}
C_{shortage} = 0
\end{equation*}

Por lo tanto, el costo total real de la solución es:

\begin{equation*}
C_{total} = 300 + 900 + 0 = 1200
\end{equation*}

\begin{table}[H]
\centering
\caption{Costo real de la solución propuesta por el operador}
\begin{tabular}{lc}
\toprule
Componente & Valor \\
\midrule
Distancia T1 & 60 km \\
Distancia T2 & 90 km \\
Distancia total & 150 km \\
Costo por distancia & \$300 \\
Costo fijo T1 & \$500 \\
Costo fijo T2 & \$400 \\
Penalización por shortage & \$0 \\
\midrule
Costo total real & \$1200 \\
\bottomrule
\end{tabular}
\end{table}

Por lo tanto, el costo declarado por el operador es incorrecto. El costo de distancia no corresponde a \$260, sino a \$300. Además, si se consideran los costos fijos de los camiones utilizados, el costo total real asciende a \$1200.

\subsection{Conclusión del Análisis}

A partir de la verificación realizada a la solución obtenida por el operador, se puede concluir que esta misma cumple con las restricciones de capacidad, satisface la demanda total sin shortage y respeta los horarios de atención, ya que los camiones esperan hasta el inicio de su ventana correspondiente.

Sin embargo, la solución presenta inconcistencias destacables. Los fill ratios declarados no coinciden realmente con las cargas reales que son requeridas, el camión $T2$ incumple de manera leve la restricción de estabilidad, ya que alcanza una diferencia de 0.3056, superando el límite permitido de 0.30, y el costo declarado está incorrecto. Considerando una distancia total de 150 km, el costo por distancia corresponde a \$300 y sumando el costo real, asciende a \$1200. Por lo tanto, la solución es parcialmente factible, pero no es completamente válida.